\section{Conclusion}
\label{sec:conclusion}

% ── Instructions ─────────────────────────────────────────────────────────
% 1 paragraph: restate problem + approach.
% 1 paragraph: key quantitative findings.
% 1 paragraph: limitations + future work.
% No new results or claims here.
% ─────────────────────────────────────────────────────────────────────────

We presented a method for reducing SLAM drift in indoor environments by
injecting camera-to-robot pose constraints into the Ceres factor graph of
a graph-based LiDAR SLAM backend.
Fixed infrastructure cameras detect the robot visually via a YOLO model,
and each detection provides an absolute pose anchor that supplements
odometry and loop-closure edges.
The approach is implemented as a standalone node requiring no
modification to the SLAM solver core.

Experiments in a 10-room ring simulation show that enabling camera
constraints reduces Absolute Trajectory Error by \textbf{?\%} and
improves map overlap with ground truth from \textbf{?\%} to
\textbf{?\%}.
These gains hold across randomised initial conditions and are robust to
constraint covariance values in the range
$\sigma_x \in [0.08, 0.15]\ \si{\metre}$.

Several directions remain for future work.
An important practical limitation of conventional infrastructure-assisted SLAM
is that camera poses are often registered only after a map has already been
created, which means any residual map drift is baked into the camera anchors
themselves. This motivates formulations in which map structure and camera
poses are refined together rather than in separate stages.
First, replacing the world-file ground-truth camera pose with a visual
relative-pose estimator (homography or PnP) will remove the dependency on
externally specified camera positions.
Second, extending the formulation to multiple robots with shared
infrastructure cameras could benefit collaborative mapping. % TODO: cite
Third, an online covariance estimator based on detection confidence would
make the method self-tuning.
